\documentclass{article}
\usepackage{amsmath}
\usepackage{amssymb}
\usepackage{bm}

\title{Fisher-Rao Geodesic: Theory and Implementation}
\author{}
\date{}

\begin{document}

\maketitle

\section{Overview}

The Fisher-Rao distance is computed using the geodesic method of Eriksen (1987) on the manifold of multivariate normal distributions equipped with the Fisher information metric (Skovgaard 1984).

\section{Canonical Parameters}

Two multivariate normal distributions $N(\boldsymbol{\mu}_1, \boldsymbol{\Sigma}_1)$ and $N(\boldsymbol{\mu}_2, \boldsymbol{\Sigma}_2)$ are transformed to canonical parameters relative to $(0, I)$ via:
\[
\boldsymbol{\mu}_{\text{trans}} = U(\boldsymbol{\mu}_2 - \boldsymbol{\mu}_1), \quad 
\boldsymbol{\Sigma}_{\text{trans}} = U\boldsymbol{\Sigma}_2 U^T
\]
where $U$ is the upper triangular Cholesky factor of $\boldsymbol{\Sigma}_1^{-1}$.

The canonical parameters are:
\[
D_e = \boldsymbol{\Sigma}_{\text{trans}}^{-1}, \quad 
\mathbf{d}_e = D_e \boldsymbol{\mu}_{\text{trans}}
\]

\section{Geodesic Computation: General Case}

\subsection{Diagonalization}

Diagonalize the precision matrix $D_e$ via eigendecomposition:
\[
D_e = V D_e^{\text{diag}} V^T
\]
where $V$ contains eigenvectors and $D_e^{\text{diag}} = \text{diag}(\lambda_1, \ldots, \lambda_p)$ contains eigenvalues.

Transform the canonical parameters:
\[
\mathbf{d}_{\text{rotated}} = V^T \mathbf{d}_e, \quad 
D = \sqrt{\text{diag}(D_e^{\text{diag}})}, \quad 
\mathbf{d} = D^{-1} \mathbf{d}_{\text{rotated}}
\]

\subsection{Block Structure of $A(y)$}

The matrix $A(y) = \log(\Lambda_0(y))$ has a block structure:
\[
A(y) = \begin{bmatrix}
B(y) & \mathbf{x}(y) & C(y) \\
\mathbf{x}(y)^T & e(y) & \mathbf{z}(y)^T \\
C(y)^T & \mathbf{z}(y) & E(y)
\end{bmatrix}
\]
where:
\begin{itemize}
  \item $B(y)$ is $p \times p$
  \item $\mathbf{x}(y)$ is $p \times 1$
  \item $C(y)$ is $p \times p$
  \item $e(y)$ is scalar
  \item $\mathbf{z}(y)$ is $p \times 1$
  \item $E(y)$ is $p \times p$
\end{itemize}

\subsection{Optimization Objective}

The geodesic is found by minimizing:
\begin{equation}
f(y) = \text{tr}(C(y)^2) + \text{tr}((B(y) + E(y))^2) + \sum(\mathbf{x}(y) + \mathbf{z}(y))^2 + e(y)^2
\label{eq:objective}
\end{equation}

At the optimum $y^*$, the matrix $A(y^*)$ satisfies the geodesic constraints:
\begin{align}
C(y^*) &= 0 \quad \text{(off-diagonal block vanishes)} \\
e(y^*) &= 0 \quad \text{(middle scalar vanishes)} \\
\mathbf{z}(y^*) &= -\mathbf{x}(y^*) \quad \text{(lower-left couples to upper-right)} \\
E(y^*) &= -B(y^*) \quad \text{(lower-right is negative of upper-left)}
\end{align}

These constraints enforce the block structure:
\[
A(y^*) = \begin{bmatrix}
B & \mathbf{x} & 0 \\
\mathbf{x}^T & 0 & -\mathbf{x}^T \\
0 & -\mathbf{x} & -B
\end{bmatrix}
\]

\subsection{Distance Computation}

The Fisher-Rao distance is computed from the optimal $A$ matrix as:
\[
D_0 = \frac{1}{2}\sqrt{\text{tr}(A(y^*)^2)}
\]

\section{Eigenvector Case}

When the canonical parameter $\mathbf{d}_e$ is an eigenvector of $D_e$ (which includes all univariate cases $p=1$), the geodesic computation simplifies to a closed-form formula.

Given:
\[
D_e \mathbf{d}_e = \lambda \mathbf{d}_e
\]

Define the mean parameter $\boldsymbol{\mu} = D_e^{-1} \mathbf{d}_e$ and $\tau = \boldsymbol{\mu}^T \mathbf{d}_e$.

Construct $\Lambda_0$ with blocks:
\begin{align}
D_e &= \text{precision matrix} \\
\mathbf{d}_e &= \text{canonical mean parameter} \\
\boldsymbol{\Phi}_0 &= -\frac{1}{2}\mathbf{d}_e \boldsymbol{\mu}^T \\
\boldsymbol{\gamma}_0 &= -\left(1 + \frac{\tau}{2}\right) \boldsymbol{\mu} \\
\boldsymbol{\Gamma}_0 &= D_e^{-1} + \left(1 + \frac{\tau}{4}\right) \boldsymbol{\mu} \boldsymbol{\mu}^T
\end{align}

Then:
\[
\Lambda_0 = \begin{bmatrix}
D_e & \mathbf{d}_e & \boldsymbol{\Phi}_0 \\
\mathbf{d}_e^T & 1 + \tau & \boldsymbol{\gamma}_0^T \\
\boldsymbol{\Phi}_0^T & \boldsymbol{\gamma}_0 & \boldsymbol{\Gamma}_0
\end{bmatrix}
\]

Compute $A = \log(\Lambda_0)$ and the distance:
\[
D_0 = \frac{1}{2}\sqrt{\text{tr}(A^2)}
\]

\section{Implementation Notes}

\subsection{Coordinate Transformations}

The implementation uses an orthogonal similarity transformation to rotate between spaces:
\[
S = \text{block-diag}(V, 1, V)
\]

The relationship between rotated and canonical coordinates is:
\[
\Lambda_{\text{canonical}} = S \Lambda_{\text{rotated}} S^T
\]

This ensures that logarithmic and exponential operations preserve the block structure across coordinate systems.

\subsection{Multi-Start Optimization}

To avoid poor local minima, the optimization is run multiple times with different random initializations. Convergence is checked and warnings are issued if $f(y)$ does not reach $< 10^{-4}$.

\subsection{Ill-Conditioning}

The computation detects ill-conditioned precision matrices ($\kappa(D_e) > 10^{10}$) and issues warnings, as these lead to numerical instability in the geodesic computation.

\section{References}

\begin{itemize}
  \item Skovgaard, L. T. (1984). A Riemannian geometry of the multivariate normal model. \textit{Scandinavian Journal of Statistics}, 11(4), 211--223.
  
  \item Eriksen, P. S. (1987). Geodesics connected with the Fisher metric on the multivariate normal manifold. Proceedings of the GST Workshop.
  
  \item Moakher, M. (2005). A differential geometric approach to the geometric mean of symmetric positive-definite matrices. \textit{SIAM Journal on Matrix Analysis and Applications}, 26(3), 735--747.
\end{itemize}

\end{document}
